\section{The terminal determinant ratio}
\label{sec:interacting-response}

The argument in this and the following sections is the terminal
response proof of \cite[\S23]{QED3Master}.
The free coefficients do not by themselves control the
reference Wick products. We must compare the normalized
interacting generating function with its one-cell free
factors, including mixed derivatives in different cells.
The quantity to be estimated is
\begin{equation}\label{eq:interacting-reference-error}
 \Theta_m^{(0)}=
 \left|\left(\prod_{c\in\mathcal C_*}
       (\partial_{t_c}^m-c_m)\mathring\Phi\right)(0)\right|^{1/N_*},
 \qquad m=2,4,\qquad N_*=|\mathcal C_*|=8L^3,
\end{equation}
where $c_m=N_0^{(m)}(0)$. The exact relation to Wick
products is \eqref{eq:external-Wick-product}.
The proof uses decay between separate insertions,
an analytic logarithm for the terminal fermion integral,
and a separate estimate of the complement of the
small-field set.

\subsection{Analyticity and normalization}
\label{subsec:interacting-determinant-ratio}

For the marked connections of
\eqref{eq:interacting-marked-connections}, set
\begin{equation}\label{eq:interacting-mark-radii}
 C_b=1+\sqrt h\,C_\beta,\qquad
 t_x=\frac{e_*^{3\varepsilon}}{C_b},\qquad
 t_x^+=\frac{e_*^{3\varepsilon}}{\sqrt h\,C_\beta},
 \qquad r=t_x/4.
\end{equation}
If $C_\beta=0$, interpret $t_x^+=\infty$.
For $t$ in the polydisc
$\mathbb D(t_x^+)=\{|t_c|<t_x^+\text{ for every }c\}$,
all the rescaled connections $a(t)^{[j]}$
belong to $\mathfrak X_j$ and
$e_*|\mathscr C(t)|_\infty<e_*^{3\varepsilon}$.
Indeed, \eqref{eq:interacting-external-probe-bound}
and $e_*\le e\sqrt h$ give
$|a(t)|_{\rm ext}<e_*^{-1+3\varepsilon}$.
This bounds the stop sup and first-difference
norms by the defining thresholds of $\mathfrak X_*$.
For $i=k_*-j$ the rescaling multiplies
the sup norm by $L^{-i/2}$ and the first
differences by $L^{-3i/2}$, while
$e_j=L^{-i/2}e_*$. Since
$\varepsilon<1/3$, the rescaled bounds are
strictly below $e_j^{-1+2\varepsilon}$ and
$e_j^{-1+3\varepsilon}$, respectively.

We first control the determinant activities accumulated
along the flow. Recall the local determinant
function $E_j^d(X,A)$ of
\cite[Proposition~6.1]{FermionBlock} and set
$E_j^{\rm det}(X,A)=E_j^d(X,A)-E_j^d(X,0)$.
Its analyticity, locality, real gauge invariance,
and evenness will all be used below; the relevant
bound is
\begin{equation}\label{eq:interacting-determinant-input}
 |E_j^d(X,A)|\le C_dM^3\mathrm e^{-(\kappa+1)d_M(X)}
                       \qquad(A\in\widetilde G_j).
\end{equation}
For a level-$(j+1)$ bond cochain $C$, write
$C_L(b)=L^{-1/2}C(b/L)$ for its one-step rescaling.
Let $\mathfrak P_0=0$ and
\begin{equation}\label{eq:interacting-P-recursion}
 \mathfrak P_{j+1}(Y,C)=
 \sum_{\overline X^{\,L}=LY}
       \bigl(\mathfrak P_j(X,C_L)+E_j^{\rm det}(X,C_L)\bigr).
\end{equation}
Then $\mathfrak P_j(X,\cdot)$ is analytic on
$\mathfrak X_j$, vanishes at zero, and depends only
on bonds meeting $X$. To see locality under
reblocking, a bond meeting $X$ has, after division
by $L$, an endpoint in $Y$ whenever
$\overline X^{\,L}=LY$. Analyticity follows from
$C_L\in\mathfrak X_j$.

For $d_M(Y)\ge1$, the reblocking estimate gives
\[
 \sum_{\overline X^{\,L}=LY}\mathrm e^{-\kappa d_M(X)}
 \le c_{\rm rb}L^3\mathrm e^{-L(\kappa-\kappa_0-1)d_M(Y)}.
\]
With
$L(\kappa-\kappa_0-1)\ge\kappa$ and
$c_{\rm rb}L^3\mathrm e^{-[L(\kappa-\kappa_0-1)-\kappa]}\le1/2$,
the weighted supremum $q_j$ of $\mathfrak P_j$
over $d_M(X)\ge1$ satisfies
$q_{j+1}\le(q_j+2C_dM^3)/2$. Since $q_0=0$,
\begin{equation}\label{eq:interacting-P-bound}
 |\mathfrak P_j(X,C)|\le2C_dM^3\mathrm e^{-\kappa d_M(X)}
                           \qquad(d_M(X)\ge1).
\end{equation}
At $j=0$ the determinant activity includes all
$192$ chamber components. The primitive bound and
reblocking estimate have the same constants at
that first step.

The determinant identity
\cite[(33)]{NonlinearStep} gives
\begin{align}
 \frac{\det(D^{\rm ch}[qa]+\bar m\eta)}{\det(D^{\rm ch}[0]+\bar m\eta)}
 &=\exp\left(\sum_X\mathfrak P_*(X,a^{[k_*]})\right)
                    \frac{\det D_*(a^{[k_*]})}{\det D_*(0)},
                           \label{eq:interacting-determinant-telescope}\\
 \sum_X\mathfrak P_*(X,a^{[k_*]})
 &=\sum_{j<k_*}\sum_XE_j^{\rm det}(X,a^{[j]}).
                           \label{eq:interacting-P-telescope}
\end{align}
Here $a^{[j]}\in\mathfrak X_j$ for every $j$.
The fine determinant uses the phase-argument convention
of \eqref{eq:response-free-ratio}: its connection argument
is $qa(t)=\sqrt\eta\sum_ct_c\widehat{\mathfrak a}_c$.
For a single cell this is precisely the phase connection
$q(t/e)\widehat a_c$ in \eqref{eq:response-free-ratio}.
Write $n_j=\ell_j^3$ for the number of unit sites at level $j$;
the spinor dimension is $4n_j$ for $j\ge1$, and $192n_0$
at the incoming chamber level. The one-step identity reads
\[
 \det D_{j+1}(a^{[j+1]})
       =b^{4n_{j+1}}\det D_j(a^{[j]})/\delta Z_j(a^{[j]}).
\]
Divide it by its zero-field version and use
\cite[(78)]{FermionBlock}. The resulting
ratio acquires the factor
\[
 \exp\left(-\sum_XE_j^{\rm det}(X,a^{[j]})\right)
\]
at each
step, which proves the first identity by iteration.
Summing \eqref{eq:interacting-P-recursion} over $Y$
proves the second: each fine polymer has exactly
one reblocked image. The zero-field determinants
are nonzero by the real covariance formula.
At $j=0$ the incoming determinant has dimension
$192n_0$, the outgoing one dimension $4n_1$;
the factor $b^{4n_1}$ cancels in the ratio.

Fix a coarse real photon field $A$ with
$\chi_*(A)=1$, and a real twist $\theta$.
The real gauge construction in
\cite[(32)]{MaxwellBlock}--%
\cite[(33)]{MaxwellBlock}
gives a representative $A^g$ with
$|A^g|_\infty\le2\ell_*^2p_*$.
The curvature bound and
\cite[Theorem~B]{MaxwellBlock}
put both $A+c_\theta$ and $A^g+c_\theta$
in $\mathrm{Dom}_{k_*}$.
Write
\[
 \mathscr A^0=H_*^{\rm ph}A^g,\qquad
 c_{\theta,\mu}=\frac{\theta_\mu}{e_*\ell_*},
 \qquad \mathscr A=\mathscr A^0+c_\theta.
\]
Summing the three photon kernel estimates in
\cite[(14)]{MaxwellBlock} bounds each
field norm of $\mathscr A^0$ by
$2C_0\ell_*^2p_*<e_*^{-\varepsilon}/2$;
the strict inequality is the corresponding
condition in \cite[(43)]{TerminalPressure}.
The local photon domains and their invariance
under real constants then give
$s\mathscr A^0+c_\theta\in\frac14\widetilde R_*$
for $0\le s\le1$. In particular all
connections needed below belong to $\mathcal W$.

The flow bound \cite[(66)]{CoverTuning}
gives an activity with even Grassmann kernels,
with equal barred and unbarred degrees, and
\begin{equation}\label{eq:interacting-activity-input}
 \|\mathscr E_*\|_*^x
       \le4e_*^{1/4-6\varepsilon},\qquad
 |m_*|\le e_*^{3/4-8\varepsilon}/8.
\end{equation}
The coefficient norm includes the weights
$\eta_*^3/48$ for a fine spinor index and
$(\eta_*^3/48)^2$ for a pair index, with all
admissible pairs, including pairs inside a cluster.
Its polymer factor is $\mathrm e^{\kappa d_M(X)}$.
For $(A,C)\in\widetilde R_*\times\mathfrak X_*$ define
\begin{equation}\label{eq:interacting-two-field-integral}
 Z_{\rm in}(A,C)=\int D\Psi\,
 \exp\left\{-\langle\bar\Psi,D_*(A+C)\Psi\rangle
       -m_*\langle\bar\psi,\psi\rangle
       +\sum_X\mathscr E_*(X;A,C,\psi,\chi)\right\},
\end{equation}
where $\psi,\bar\psi$ are substituted using
$H_*^\#(A+C)\Psi$ and their barred partners,
and $\chi,\bar\chi$ using the corresponding
covariant Hölder quotients.
The inclusion
$\widetilde R_*+\mathfrak X_*\subset\widetilde G_*$
makes all these maps analytic; the activity
kernels are analytic in $(A,C)$.
Berezin integration is coefficient extraction,
and the exponential is its scalar exponential
times a finite polynomial in the positive-degree
coefficients. Thus $Z_{\rm in}$ is analytic
on this two-field domain.
Set
\begin{align}
 Z_{\rm tm}(t)&=Z_{\rm in}(\mathscr A,\mathscr C(t)),
                               \label{eq:interacting-terminal-integral}\\
 f_{\rm int}(t)&=
   \frac{Z_{\rm tm}(t)}{\det D_*(c_\theta+\mathscr C(t))}.
                               \label{eq:interacting-factor-definition}
\end{align}
The Berezin measure has one pair for each
of the $4\ell_*^3$ coarse spinor coordinates.
The two-field analyticity gives analyticity of
$Z_{\rm tm}$ on
$\mathbb D(t_x^+)$.

\begin{proposition}\label{prop:interacting-factor}
Suppose
\begin{gather}
 4C_0\ell_*^3g_1\le e_*^{-1/4+\varepsilon},\qquad
 4\ell_*^3(2C_K+K_0+4)\le e_*^{-2\varepsilon},\qquad
 e_*^{1/4-9\varepsilon}\le1/8,
                                      \label{eq:interacting-factor-conditions}\\
 48e_*^{2\varepsilon}\le\mu_- .
                                      \label{eq:interacting-small-complex-mass}
\end{gather}
Then $f_{\rm int}$ is analytic and
\begin{equation}\label{eq:interacting-factor-bound}
 |f_{\rm int}(t)-1|\le\frac12e_*^{1/4-9\varepsilon}
                              \le\frac1{16}.
\end{equation}
In particular $Z_{\rm tm}$ has no zero on the
polydisc. The normalized terminal ratio
\begin{equation}\label{eq:interacting-terminal-ratio}
 R_{\rm tm}(t;\theta,A)=
 \frac{Y_*^{\rm x}(A+c_\theta;\mathscr C(t))}
      {Y_*^{\rm x}(A+c_\theta;0)}
 =\exp\left(\sum_X\mathfrak P_*(X,\mathscr C(t))\right)
       \frac{Z_{\rm tm}(t)}{Z_{\rm tm}(0)}
\end{equation}
is nonzero, continuous in $(t,\theta,A)$ on
$\mathbb D(t_x^+)\times\R^3\times\{\chi_*=1\}$,
and analytic in $t$. It has a unique continuous
logarithm $\mathcal G(t;\theta,A)$ that vanishes
at zero, and this logarithm is analytic in $t$.
\end{proposition}

\begin{proof}
Use the covariance
\[
 C_t=D_*(c_\theta+\mathscr C(t))^{-1}.
\]
The inverse exists by
\cite[(75)]{TerminalPressure}, because
$e_*|\mathscr C(t)|_\infty<e_*^{3\varepsilon}
<\mu_*/24$. Each covariance row has sum at
most $4\ell_*^3g_1$. Put
$a_w=C_0^{-1}e_*^{-1/4+\varepsilon}$.
The first condition in
\eqref{eq:interacting-factor-conditions} gives
$a_w\ge4\ell_*^3g_1\ge1$.

Separate the quadratic form with covariance $C_t$
from the remaining even interaction $V_t$.
Normalized Gaussian Grassmann integration gives
\begin{equation}\label{eq:interacting-crude-Gaussian}
 |f_{\rm int}(t)-1|\le \mathrm e^{|V_t|_{a_w}}-1.
\end{equation}
To justify this particular norm bound, replace
the unbarred generators by $C_t\Psi$ and leave
the barred generators fixed. Their row sums
are at most $a_w$, so this substitution maps
the $a_w$ coefficient norm into the unit
coefficient norm. Integration against the identity
covariance has norm one in that latter norm.
Apply this to $\mathrm e^{V_t}-1$, and use the
submultiplicativity of the coefficient norm.

There are three terms in $V_t$. The difference
between the two coarse quadratic forms has
entry sum at most
$C_Ke_*^{1-\varepsilon}\ell_*^3/2$, by
\cite[(66)]{TerminalPressure} integrated along
$c_\theta+\mathscr C(t)+s\mathscr A^0$.
A quadratic Grassmann form has coefficient norm
at most $4a_w^2$ times this block sum, so
its contribution is bounded by
$2C_K\ell_*^3e_*^{1/2+\varepsilon}$.
The terminal mass estimate and
\eqref{eq:interacting-activity-input} give respectively
\[
 \frac12\ell_*^3e_*^{1/4-6\varepsilon},
 \qquad K_0\ell_*^3e_*^{1/4-7\varepsilon}.
\]
For the activity sum, substitute the terminal
field maps in the coefficient norm, then sum
$\mathrm e^{-\kappa d_M(X)}$ by assigning each polymer
one of its $M$-cubes and using the polymer
sum bound $K_0$. We have used
$4e_*^\varepsilon\le1$.
Consequently
\[
 |V_t|_{a_w}\le
 \ell_*^3(2C_K+K_0+4)e_*^{1/4-7\varepsilon}
 \le\frac14e_*^{1/4-9\varepsilon}.
\]
The inequality $\mathrm e^x-1\le2x$ proves
\eqref{eq:interacting-factor-bound}.
The denominator in \eqref{eq:interacting-factor-definition}
is analytic and nonzero, so $f_{\rm int}$ is analytic.

By \cite[(16)]{TerminalPressure} the terminal density has the form
\[
 Y_*^{\rm x}(A';C)=
 (\det D_*(0))^{-1}
 \exp\left(-\varepsilon_*^{\rm v}\ell_*^3+
                         \sum_X\mathfrak P_*(X,C)\right)
 Z_{\rm in}(H_*^{\rm ph}A',C).
\]
Its gauge invariance in \cite[Theorem~B]{TerminalPressure}
permits replacing $A$ by
$A^g$. The photon minimizer preserves real
constants. Dividing the density at $C=\mathscr C(t)$
by its value at zero proves
\eqref{eq:interacting-terminal-ratio}; the quotient
also shows that the result is independent of
the chosen gauge representative. For continuity,
use the ungauged linear map
$A+c_\theta\mapsto H_*^{\rm ph}(A+c_\theta)$
in the displayed formula. The terminal Grassmann
integral is analytic in its two field arguments,
and its zero-mark denominator is nonzero.

The required logarithm is
\begin{equation}\label{eq:interacting-logarithm}
 \mathcal G(t;\theta,A)=
 \int_0^1\sum_ct_c
       \frac{\partial_{t_c}R_{\rm tm}(st;\theta,A)}
            {R_{\rm tm}(st;\theta,A)}\,\dd s .
\end{equation}
Local Cauchy formulas show that the derivatives
in this integrand are jointly continuous;
integration over $[0,1]$ preserves continuity
and analyticity. Along a ray the integrand is
the logarithmic derivative of $R_{\rm tm}(st)$.
Differentiating the product of this ratio with
the exponential of minus the integrated
logarithmic derivative gives zero, and its
value at $s=0$ is one. Hence
$\mathrm e^{\mathcal G}=R_{\rm tm}$.
Any two continuous logarithms differ by
$2\pi i$ times an integer-valued continuous
function on the connected polydisc. Their
common value at zero proves uniqueness.
\end{proof}

